% !TeX root = main.tex

\chapter{Introduction}

Computing the distance between two points is a common geometrical task.
A natural extension to that would be to find the \textit{closest pair} of points,
i.e., the pair that yields the smallest distance. Assuming a point in $d$-dimensional
space, computing the distance between two points takes $\mathcal{O}(d)$ time.
There are $\binom{n}{2}$ pairs to check, forcing a total time complexity of
$\mathcal{O}(d \cdot n^{2})$.

In general, for large $d$, it's not possible to beat quadratic time in the \textit{worst case}, assuming the
\textbf{Strong Exponential Time Hypothesis (SETH)}. Luckily for us, there are tricks to salvage
the situation at low dimensions, and for practical applications like vector databases,
people have ditched the intractable task of obtaining an exact solution, opting for
approximate algorithms instead.

\section{Duplicate Points}

Depending on the circumstances in which these points were generated, it's quite possible that we
end up with duplicate points. For example, consider the Birthday Paradox. While it seems strange
that only $23$ people are required to produce a coin flip chance of having a shared birthday
(the equivalent of two overlapping points), that's just how the probability checks out.

One way to solve this problem would be to first sort the points and check adjacent items to see if
they're equal. Another approach is to inject them into a \textit{hash set} one by one, and prior to the
insertion, if the set already contains the point, we know a duplicate exists. The extra logarithmic
factor in the former might seem burdensome, but it might actually beat the latter in terms of real
runtime, considering hash sets in many languages are often plagued by high constant factors. Not to
mention sorting also happens to be the first step in many of the subsequent algorithms, so the duplicate
check arrives almost free in a sense.

\section{Points on a Line ($d=1$)}

At $d=1$, the points are just an array $a$ of scalar values, so one may simply sort $a$ and
compare adjacent elements, consuming
$\mathcal{O}(n \log n)$ time. If the range of values ($k = \max(a) - \min(a)$)
is small enough, \textit{counting sort} may be exploited to solve the problem
in $\mathcal{O}(n + k)$.

At this juncture, it's natural to question if it's possible to achieve $\mathcal{O}(n)$
irrespective of $k$, and as it turns out, there are randomised
algorithms with linear expected time. These algorithms are rooted in
the idea of \textit{decomposition}, and they generalise to higher dimensions as well.
We'll elaborate more on the details when we get there.
